\documentclass[11pt,a4paper]{article}
\usepackage[utf8]{inputenc}
\usepackage[T1]{fontenc}
\usepackage{amsmath,amssymb,bm,array,geometry,color,fancyhdr,hyperref}
\geometry{left=19mm,right=19mm,top=18mm,bottom=19mm}
\definecolor{ink}{RGB}{20,52,72}
\definecolor{accent}{RGB}{22,115,139}
\hypersetup{colorlinks=true,urlcolor=accent,linkcolor=ink}
\setlength{\parindent}{0pt}\setlength{\parskip}{3pt}
\setlength{\abovedisplayskip}{7pt}\setlength{\belowdisplayskip}{7pt}\setlength{\jot}{4pt}
\pagestyle{fancy}\fancyhf{}\renewcommand{\headrulewidth}{0pt}
\fancyfoot[L]{\small\color{ink}pysnspd / modelo efectivo final / 12 de octubre de 2026}
\fancyfoot[R]{\small\thepage}
\newcommand{\block}[1]{\par\vspace{5pt}{\large\bfseries\color{ink}#1}\par\vspace{3pt}}
\newcommand{\labelrow}[2]{\(\displaystyle #1\)& #2\\}
\newcommand{\sourcecode}[1]{\par{\footnotesize\color{ink}\textbf{Implementación:} \texttt{#1}}\par}
\newcommand{\origin}[1]{\par{\footnotesize\color{ink}\textbf{Campo / procedencia:} #1}\par}
\newcommand{\dd}{\,\mathrm d}
\newcommand{\Eu}{E_u}
\newcommand{\Tb}{T_b}
\newcommand{\Rea}{\operatorname{Re}}
\newcommand{\Ima}{\operatorname{Im}}
\newcommand{\FD}{f_{\rm FD}}
\newcommand{\BE}{n_{\rm BE}}
\newcommand{\Uscale}{u_i^*}
\newcommand{\Rcut}{P_{\rm res}}
\newcommand{\Fgap}{\mathcal F_i}
\newcommand{\meV}{\,\mathrm{meV}}
\begin{document}
{\LARGE\bfseries\color{ink}Cantidades precomputadas}\par
{\large Escalas, cuadraturas, material y preparación inicial}\par
{\small Entradas exploratorias explícitas. Fuentes externas preservadas; ningún cálculo Quantum ESPRESSO/DFPT propio.}
\renewcommand{\theequation}{3.\arabic{equation}}
\block{Escalas SI y malla Delaunay--Voronoi}
\begin{gather}
 E_u=k_BT_c,\qquad
 \ell_0=\sqrt{\frac{\hbar D}{2k_BT_c}},\qquad
 \tau_K=\frac{\pi\hbar}{8k_BT_c},\\
 \sigma_n=\frac1{R_\square d_f}=2e^2N_0D,\qquad
 N_0=\frac1{2e^2DR_\square d_f},\qquad
 (R_\square d_f)_{\mathrm{mismo}\ N_0,D}=\mathrm{constante},\\
 m_i=A_i^{dual}/\ell_0^2,\qquad
 c_{ij}=L_{ij}^{dual}/L_{ij},\qquad V_i=d_fA_i^{dual},\\
 A_{ij}^{face}=d_fL_{ij}^{dual},\qquad
 {\cal D}_{ij}=D A_{ij}^{face}/L_{ij},\qquad
 u_i^*=N_0E_u^2V_i,\\
 T_c=8.65\,\mathrm K,\quad T_b=0.9\,\mathrm K,\quad
 D=5\times10^{-5}\,\mathrm{m^2/s},\quad R_\square=608\,\Omega,\\
 d_f=7\,\mathrm{nm},\qquad W=80\,\mathrm{nm},\qquad
 h_{\rm mesh}\in\{5,3.5\}\,\mathrm{nm},\quad
 L\in\{160,320,480\}\,\mathrm{nm},\\
 N_0=9.1532978211\times10^{46}\,\mathrm{J^{-1}m^{-3}},
 \qquad \ell_0\simeq4.698490897\,\mathrm{nm}.
\end{gather}
\block{Calibradores físicos y nomenclatura Allmaras}
\begin{gather}
 \tau_{ee}(T_c)>0,\quad\tau_{ep}(T_c)>0,\qquad
 \tau_{ee}(T_c)_{\rm nominal}=5\,\mathrm{ps},\quad
 \tau_{ep}(T_c)_{\rm nominal}=24.7\,\mathrm{ps},\\
 \tau_{\rm esc}^{\rm nominal}=15\,\mathrm{ps}.
\end{gather}
{\small \(\tau_{ee}(T_c),\tau_{ep}(T_c)\): parámetros libres, pendientes de calibración. Valores nominales de referencia; no centro ni límites de un rango. \(\tau_{\rm esc}\): hipótesis de interfaz separada. \(h_{\rm mesh}\) es separación de malla; \(d_f\) es espesor físico.}
\par\noindent\begin{tabular}{@{}p{.30\linewidth}p{.65\linewidth}@{}}
\hline
\labelrow{A_i^{dual},L_{ij}^{dual}}{área de Voronoi [m$^2$]; longitud de cara dual [m]}
\labelrow{L_{ij},m_i,c_{ij}}{arista Delaunay [m]; área y conductancia normalizadas}
\labelrow{W,L,d_f,h_{\rm mesh}}{ancho, longitud resuelta, espesor y resolución [m]}
\labelrow{\sigma_n,N_0,D}{conductividad [S/m]; DOS monoespín; difusión [m$^2$/s]}
\labelrow{\tau_K}{escala cinética KWT [s]}
\labelrow{\tau_{ee}(T_c),\tau_{ep}(T_c)}{calibradores libres a $T_c$ [s]}
\labelrow{\tau_{\rm esc}}{tiempo de escape fonónico [s]}
\hline
\end{tabular}\par
\vfill
\origin{Relación de Einstein, geometría de volúmenes finitos y KWT. \href{https://thesis.caltech.edu/13748/}{Allmaras (2020)}, ec.\ 3.24, pp.\ 94--98; escenario experimental motivado por Korzh (2020), sin afirmar calibración de su película.}
\sourcecode{effective\_bcs\_preparation.py; effective\_coupled.py}

\newpage
\block{Estados nativos fijos y emparejamiento renormalizado}
\begin{gather}
 \xi_k=E_u\bar\xi_k>0,\qquad
 \int_0^{\Lambda/E_u}F(\bar\xi)\dd\bar\xi
 \simeq\sum_kw_kF(\bar\xi_k),\qquad w_k>0,\\
 \frac1g=\sum_kw_k\frac{\tanh(\bar\xi_k/2)}{\bar\xi_k},\qquad
 \left[\frac1g-\sum_k\frac{w_k[1-2\FD(E_u\bar\xi_k,T_c)]}{\bar\xi_k}\right]=0,\\
 [0,\Lambda/E_u]=\bigcup_b[a_b,a_{b+1}],\quad
 \bar\xi_{bk}=\frac{a_b+a_{b+1}}2+\frac{a_{b+1}-a_b}2x_k,\\
 w_{bk}=\frac{a_{b+1}-a_b}2w_k^{GL},\qquad
 \{a_b\}=\{0,\tfrac18,\tfrac14,\tfrac12,1,2,4,8,16,32,\ldots\},\\
 \Lambda/E_u\in\{32,64\},\qquad
 n_{GL}\in\{6,8\}.
\end{gather}
\block{Cuadratura fonónica y regla fija de caras}
\begin{gather}
 0<\omega_0<\cdots<\omega_{Q-1}=\Omega_D/E_u,\qquad
 W_q^{ph}=E_u\bar W_q^{ph},\\
 \bar W_0^{ph}=\frac{\omega_1-\omega_0}2,\quad
 \bar W_q^{ph}=\frac{\omega_{q+1}-\omega_{q-1}}2,\quad
 \bar W_{Q-1}^{ph}=\frac{\omega_{Q-1}-\omega_{Q-2}}2,\\
 c_q=\frac{g_{ph}(E_u\omega_q)W_q^{ph}}{N_0E_u},
 \qquad Q\in\{96,144\},\qquad \omega_0=10^{-8}.
\end{gather}
{\small Nodos nativos, modos fonónicos y regla de referencia \(x_a,w_a^{GL}\) precomputados. Las energías \(E_{ik}(d_i)\), los tiempos \(\tau(T)\) y el soporte, nodos y pesos transformados de cada cara se evalúan durante la evolución y se documentan en las ecuaciones principales y auxiliares. La misma cuadratura electrónica se usa en \(g\), energía y fuerza: no tabla de núcleo singular ni \(\delta_{\rm core}\).}
\par\noindent\begin{tabular}{@{}p{.30\linewidth}p{.65\linewidth}@{}}
\hline
\labelrow{\xi,\bar\xi,w_k}{energía normal BCS [J], su coordenada/$E_u$, peso adimensional}
\labelrow{\Lambda,g}{corte electrónico [J]; emparejamiento BCS renormalizado}
\labelrow{x_a,w_a^{GL},n_{GL}}{nodos y pesos de referencia fijos; orden de Gauss--Legendre}
\labelrow{Q,K,\omega_q}{número de modos fonónicos, estados electrónicos y energía/$E_u$}
\labelrow{W_q^{ph},c_q}{peso físico fonónico [J]; capacidad nativa}
\hline
\end{tabular}\par
\vfill
\origin{Teoría BCS de emparejamiento, renormalización a \(T_c\) y cuadratura positiva. Discretización nativa propia; no consulta Usadel durante la evolución.}
\sourcecode{effective\_bcs\_functional.py; effective\_bcs\_kinetics.py}

\newpage
\block{Material acústico efectivo y enlace del calibrador}
\begin{gather}
 \Omega_D=k_B\Theta_D,\qquad
 g_{ph}(\Omega)=\frac{3\Omega^2}{2\pi^2\hbar^3v_{\rm avg}^3}
 \mathbf1_{(0,\Omega_D]}(\Omega),\\
 \tau_0^{e-ph}=\frac{720\zeta(5)}{\pi^2}\tau_{ep}(T_c),\qquad
 \alpha^2F(\Omega)=\frac{\hbar\Omega^2}{\pi\tau_0^{e-ph}(k_BT_c)^3}
 \mathbf1_{(0,\Omega_D]}(\Omega),\\
 \alpha^2(\Omega)=\frac{\alpha^2F(\Omega)}{g_{ph}(\Omega)},\qquad
 \frac{2\pi}{\hbar}\alpha^2F(\Omega)\dd\Omega
 =\frac{2\omega^2}{\tau_0^{e-ph}}\dd\omega,\quad\omega=\Omega/E_u,\\
 C_{e,N}(T)=\frac{2\pi^2}3N_0k_B^2T,\qquad
 P_{e-ph,N}=\frac{96\zeta(5)N_0k_B^2}{\tau_0^{e-ph}T_c^3}
 (T_e^5-T_{ph}^5),\\
 \frac{C_{e,N}(T_c)}
 {\left.\partial_{T_e}P_{e-ph,N}\right|_{T_e=T_{ph}=T_c}}
 =\frac{\pi^2\tau_0^{e-ph}}{720\zeta(5)}=\tau_{ep}(T_c),\\
 u_{ph}(T)=\int_0^{\Omega_D}
 \frac{\Omega g_{ph}(\Omega)}{\exp[\Omega/(k_BT)]-1}\dd\Omega,\\
 C_{ph}(T)=k_B\int_0^{\Omega_D}g_{ph}(\Omega)
 \frac{x^2e^x}{(e^x-1)^2}\dd\Omega,\quad x=\Omega/(k_BT),\\
 C_{ph}(T)\simeq B_{ph}T^3,\qquad
 B_{ph}=\frac{2\pi^2k_B^4}{5\hbar^3v_{\rm avg}^3},\qquad
 u_{ph}(T)\simeq B_{ph}T^4/4,\\
 \gamma_{\rm ac}\equiv\frac{8\pi^2}{5}
 \frac{C_{e,N}(T_c)}{C_{ph}(T_c)}\simeq23.4774459,\qquad
 \frac{C_{e,N}(T_c)}{C_{ph}(T_c)}\simeq1.486726629,\\
 v_{\rm avg}=4912\,\mathrm{m/s},\qquad\Theta_D=664\,\mathrm K,\qquad
 \Omega_D\simeq57.2191\,\mathrm{meV},\\
 \tau_{ep}(T_c)_{\rm nominal}=24.7\,\mathrm{ps}
 \ \Longrightarrow\ \tau_0^{e-ph}\simeq1.868\,\mathrm{ns}.
\end{gather}
{\small Material acústico 3D efectivo, sin ajuste absoluto de la película. Cambiar \(\tau_{ep}(T_c)\) cambia de forma vinculada la movilidad KWT y \(\alpha^2F\). \(\tau_0^{e-ph}\neq\tau_K\): escalas de colisiones y del condensado, respectivamente. \(\gamma_{\rm ac}\) es un diagnóstico derivado de las capacidades actuales con \(R_\square=608\,\Omega\), no un calibrador independiente; no se importa \(\gamma=60\) de otra parametrización.}
\par\noindent\begin{tabular}{@{}p{.30\linewidth}p{.65\linewidth}@{}}
\hline
\labelrow{\alpha^2F,\alpha^2,g_{ph}}{Eliashberg adimensional; cociente [J\,m$^3$]; DOS [J$^{-1}$m$^{-3}$]}
\labelrow{v_{\rm avg},\Theta_D,\zeta}{velocidad media [m/s]; temperatura Debye [K]; zeta de Riemann}
\labelrow{u_{ph},C_{e,N},C_{ph}}{energía fonónica [J/m$^3$]; capacidades volumétricas normal electrónica y fonónica [J/(m$^3$K)]}
\labelrow{\gamma_{\rm ac},P_{e-ph,N}}{razón normalizada de capacidades a \(T_c\), adimensional; potencia normal [W/m$^3$]}
\hline
\end{tabular}\par
\vfill
\origin{Debye y cinética acústica: \href{https://thesis.caltech.edu/13748/}{Allmaras (2020)}, ecs.\ 2.6--2.8 (definición de \(\gamma_{\rm ac}\)) y p.\ 81, relación explícita \(\tau_0/\tau_{ep}(T_c)\). Normalización conservativa de Vodolazov, ecs.\ 3--4; no se adopta el factor de dispersión fonónica omitido en la ec.\ 2.7 impresa de la tesis.}
\sourcecode{effective\_material.py; effective\_bcs\_kinetics.py}

\newpage
\block{Tabla externa NbN: forma QNN/MIT preservada}
\begin{gather}
 {\mathsf D}_{RAW}=\{(\nu_j,A_j,F_j)\}_{j=1}^{11999},\qquad
 {\mathsf D}_{forma}=\operatorname{stable-first}
 [\operatorname{sort}_{\nu}({\mathsf D}_{RAW})],\quad N_{forma}=11995,\\
 F'_j=\begin{cases}F_j,&F_j>0,\\0,&F_j\leq0,\end{cases}\qquad
 A'_j=\begin{cases}A_j,&F'_j>0\ \land\ \nu_j>0,\\0,&\text{otro caso},\end{cases}\\
 \Omega_j=2\pi\hbar\,\nu_j\,10^{12}\,\mathrm{Hz/THz},\qquad
 A(\Omega)=\operatorname{interp}_{lin}(\Omega_j,A'_j),\\
 F(\Omega)=\operatorname{interp}_{lin}(\Omega_j,F'_j),\qquad
 \lambda_{e-ph}=2\int_0^\infty\frac{\alpha^2F(\Omega)}{\Omega}\dd\Omega,\\
 C_{forma}(T)=k_B\int F(\Omega)\frac{x^2e^x}{(e^x-1)^2}\dd\Omega,\qquad
 x=\Omega/(k_BT),\\
 s_F=\frac{B_{ph}T_{\rm ref}^3}{C_{forma}(T_{\rm ref})},\qquad
 g_{ph}^{forma}(\Omega)=s_FF(\Omega),\qquad T_{\rm ref}=T_c.
\end{gather}
\par\noindent\begin{tabular}{@{}p{.30\linewidth}p{.65\linewidth}@{}}
\hline
\labelrow{\nu_j}{frecuencia ordinaria THz: hipótesis explícita del lector upstream}
\labelrow{A_j,F_j}{columnas $\alpha^2F$ y DOS de forma; normalización original incompleta}
\labelrow{\lambda_{e-ph},s_F}{constante Eliashberg; escala térmica de comparación}
\labelrow{g_{ph}^{forma}}{otra referencia efectiva anclada; no se usa en el cierre acústico actual}
\hline
\end{tabular}\par
\vspace{8pt}
{\small
\textbf{Archivo:} \texttt{nbn-a2f-ph.dat}.\\
\textbf{Procedencia:} QNN / Alejandro Simon (2025), repositorio externo
\href{https://github.com/qnngroup/proj-KE-solver}{\texttt{qnngroup/proj-KE-solver}}.\\
\textbf{Revisión:} \texttt{5b6bd747f80016da5ccd51db73c110a8ecc6abf6}.\\
\textbf{SHA-256 RAW:} \texttt{e94f11273e6c42b7c8fad978b78d11916}\\
\hspace*{28mm}\texttt{6e241bec1614ee77e4176906fa8f5df}.\\
\textbf{Licencia upstream:} MIT.\\
\textbf{Unidad absoluta y base atómica de \(F_j\):} no recuperadas.\\
\textbf{\(A_j\) adimensional:} convención del lector, no metadato recuperado.\\
\textbf{Quantum ESPRESSO / DFPT propios:} \(0\).\\
\textbf{Uso actual del RAW en las tasas acústicas:} \(0\).
}
\vfill
\origin{Datos externos de espectros fonónicos/Eliashberg. Su conservación documental no acredita normalización absoluta ni reemplaza el material efectivo autorizado.}
\sourcecode{material\_preprocessing.py; effective\_material.py}

\newpage
\block{Nuevo equilibrio DC efectivo e inductancia exterior}
\begin{gather}
 {\cal F}_i(d_{\rm DC},p_{\rm DC},\alpha=0)=0\quad(i\text{ interior}),\qquad
 \sum_{i\in L}(BI^s)_i=I_{\rm DC},\\
 p_{{\rm DC},ik}=\FD(E_{{\rm DC},ik},T_b),\qquad
 n_{{\rm DC},iq}=\BE(\Omega_q,T_b),\qquad
 I_{\rm DC}=21.2972061449\,\mathrm{\mu A},\\
 a=|\Delta|\quad\text{[J]},\qquad j_s(a,q)=\frac{4e}{\hbar}K(a;T_b)a^2q,\qquad
 \ell_{k,bulk}^{diff}=\frac{\hbar}{2e}
 \left.\frac{\dd q}{\dd I}\right|_{I_{\rm DC}},\\
 L_{k,\rm ext}^{ref}=96\,\mathrm{nH}
 +(5\,\mathrm{\mu m}-L)\ell_{k,bulk}^{diff},\\
 I_b^0=I_s^0=I_{\rm DC},\quad v_c^0=0,\quad V_b=R_bI_{\rm DC},\\
 R_b=10\,\mathrm{k\Omega},\quad L_b=1\,\mathrm{\mu H},\quad
 R_l=50\,\Omega,\quad C_c=100\,\mathrm{pF}.
\end{gather}
\block{Fotón: depósito inicial único en el interior}
\begin{gather}
 E_\gamma=\frac{2\pi\hbar c}{\lambda},\quad
 E_{\rm ret}=\eta_{\rm ret}E_\gamma,\quad \eta_{\rm ret}=0.667,\quad
 \lambda\in\{775,1550\}\,\mathrm{nm},\\
 \sigma_{\rm ref}=\sqrt{2D\tau_{\rm th,ref}}=7.549834435\,\mathrm{nm},
 \qquad\tau_{\rm th,ref}=0.57\,\mathrm{ps},\\
 \sigma\in\{\sigma_{\rm ref}/\sqrt2,\ \sigma_{\rm ref},\
 \sqrt2\,\sigma_{\rm ref}\},\qquad
 \mathbf r_0=((x_{\min}+x_{\max})/2,\ y_0),\\
 a_i=\mathbf1_{i\notin L\cup R}
 \exp[-|\mathbf r_i-\mathbf r_0|^2/(2\sigma^2)],\quad
 E_{{\rm ret},i}=\frac{V_ia_i}{\sum_jV_ja_j}E_{\rm ret},\\
 u_i^{pop}(T)=4u_i^*\sum_kw_ke_{{\rm DC},ik}\FD(E_{{\rm DC},ik},T)
 +u_i^*\sum_qc_q\omega_q\BE(\Omega_q,T),\\
 E_{{\rm ret},i}=u_i^{pop}(T_i^0)-u_i^{pop}(T_b),\\
 \sum_iE_{{\rm ret},i}=E_{\rm ret},\quad
 p_{ik}^0=\FD(E_{{\rm DC},ik},T_i^0),\quad n_{iq}^0=\BE(\Omega_q,T_i^0),\\
 d_i^0=d_{{\rm DC},i},\qquad
 (I_b,I_s,v_c)^0=(I_{\rm DC},I_{\rm DC},0).
\end{gather}
{\small \(t_0=0\): transferencia efectiva postcascada; \(\tau_{\rm lat}=t_\times-t_0\), sin sumar \(0.57\) ps. Esa referencia sólo determina \(\sigma_{\rm ref}\), sin cascada óptica resuelta. DC preparado con esta energía y corriente; la admisión Usadel histórica no se reutiliza. \(L_{k,\rm ext}\) queda congelada durante el pulso. FD/BE común es hipótesis inicial tras cascada, no cierre térmico posterior. Retención y anchura son exploratorias sin calibración; \(0.57\) ps no añade una ventana térmica durante la evolución.}
\par\noindent\begin{tabular}{@{}p{.30\linewidth}p{.65\linewidth}@{}}
\hline
\labelrow{a=|\Delta|}{módulo físico del parámetro de orden [J]}
\labelrow{I_{\rm DC},q,j_s}{corriente [A], gradiente de fase [m$^{-1}$], densidad de corriente [A/m$^2$]}
\labelrow{\ell_{k,bulk}^{diff}}{inductancia cinética diferencial por longitud [H/m]}
\labelrow{E_\gamma,E_{\rm ret},\eta_{\rm ret}}{energías del fotón y retenida [J]; fracción retenida}
\labelrow{\lambda,\sigma,\mathbf r_0}{longitud de onda, anchura gaussiana y absorción [m]}
\labelrow{\tau_{\rm th,ref},T_i^0}{referencia de cascada [s]; temperatura inicial equivalente [K]}
\labelrow{u_i^{pop},c,k_B}{energía nodal [J]; velocidad de luz [m/s]; constante de Boltzmann [J/K]}
\hline
\end{tabular}\par
\vfill
\origin{Preparación DC variacional propia; circuito de la memoria. Retención motivada por Korzh (2020); escala de cascada de \href{https://arxiv.org/pdf/1611.06060}{Vodolazov (2017)} a 1.3 eV, extrapolación exploratoria autorizada para ambos fotones.}
\sourcecode{effective\_bcs\_preparation.py; effective\_bcs\_ports.py}
\vspace{4pt}
{\footnotesize
\href{https://thesis.caltech.edu/13748/}{Allmaras, tesis Caltech (2020)};\quad
\href{https://arxiv.org/pdf/1611.06060}{Vodolazov, Phys.\ Rev.\ Applied 7, 034014 (2017)};\quad
\href{https://doi.org/10.1103/PhysRev.108.1175}{Bardeen--Cooper--Schrieffer (1957)}.
}
\end{document}
